\documentclass[../main.tex]{subfiles}
\begin{document}
%==========================================TIÊU ĐỀ TẬP========================================
\begin{table}[h]
  \begin{adjustwidth}{-1cm}{}
    \begin{tabular}{>{\centering\arraybackslash}p{6cm}|>{\raggedright\arraybackslash}p{12.5cm}}
      \multirow{5}{6cm}
      % Cột bên trái: ÉPISODE và số tập
      {\\[-40pt]
        \flaregothic\fontsize{20pt}{20pt}\selectfont \centering
        \color{eptype}ÉPISODE\color{black}\\
        \burbank\fontsize{72pt}{72pt}\selectfont
        \color{epnum}84
        \\[18pt]
        \flaregothic\fontsize{14pt}{14pt}\selectfont
        \color{epattr}MISE À JOUR DU MODÈLE pour \uline{\color{Robopon} Roboco}
        % \flaregothic\fontsize{18pt}{18pt}\selectfont
        % \color{epattr}ÉPISODE PERDU
        \color{black}
      }
       &
      % Dòng thứ nhất: tên tập tiếng Nhật
      {\fotskip\fontsize{18pt}{18pt}\selectfont \ruby{体}{からだ}が\ruby{動}{うご}かない！\ruby{助}{たす}けて！}
      \\[6pt]
      % Dòng thứ hai: tên tập tiếng Anh
       & {\cafeteria\fontsize{18pt}{18pt}\selectfont A Malfunctioning Experience}                  \\[4pt]
      % Dòng thứ ba: tên tập tiếng Việt
       & {\vnmsans\fontsize{18pt}{18pt}\selectfont Hỏng!}                                          \\[8pt]
      % Dòng thứ tư: Nhân vật xuất hiện
       & {\brian{\fontsize{14pt}{14pt}\selectfont Track used: Steam Ages - Demonstration Mini-Game}
         }
      \\[8pt]
      % Dòng cuối cùng: Thời gian diễn ra của tập (thường sẽ là ngày phát hành)
       & {\continuum\fontsize{16pt}{16pt}\selectfont Ngày phát hành: 13/12/2020}                   \\[18pt]
    \end{tabular}
  \end{adjustwidth}
\end{table}
%=========================================TRUYỆN CHÍNH========================================
\color{black}

Hôm nay, Kuon có việc riêng ở nhà nên đã xin nghỉ một ngày. Văn phòng \hll\ vì thế vắng đi một bóng người quen thuộc, để lại các cô gái tự xoay xở với công việc của mình. Một ngày tưởng chừng như bình thường\dots\ cho đến khi Choco bước vào phòng và chứng kiến một cảnh tượng kỳ lạ.

Trước mặt cô, Roboco đứng bất động trong một tư thế khó hiểu: một tay giơ lên cao, tay còn lại hạ xuống, chân phải bước lên trước nhưng cả hai bàn chân lại kiễng lên như đang giữ thăng bằng trên một sợi dây vô hình. Điều đáng nói là cô nàng Rô-bốt cứ đứng yên như thế, không hề cử động dù chỉ một chút, giống như ai đó đã bấm nút ``tạm dừng'' trên một bộ phim.

\img{7-2a}

Choco chớp mắt vài lần, cố gắng lý giải tình huống trước mắt nhưng không tài nào hiểu nổi. Cô chưa kịp lên tiếng thì ngay khi cánh cửa phòng bật mở, nàng Người máy vừa ``yoo-hoo'' một tiếng. Bất ngờ thay, Roboco hơi giật nhẹ - chính xác hơn là chỉ có phần đầu cô khẽ nghiêng một chút, như thể nhận ra có người vừa vào phòng.

Nàng Y tá lập tức nhận ra điều bất thường. ``Chị đang làm gì vậy?'' cô hỏi, tiến lại gần hơn để quan sát.

``Chị bị kẹt,'' giọng Roboco vang lên một cách lạ lẫm: hơi đơn điệu, thiếu đi sự linh hoạt thường ngày và có chút méo mó như một chiếc loa bị rè.

Choco nhướng mày: ``Tại sao?''

``Dạo gần đây chị lười bảo trì quá,'' Roboco thở dài, nhưng với tình trạng hiện tại, âm thanh phát ra từ cô nghe cứ như một con robot thực thụ. ``Bây giờ ấy'' - cô cố gắng quay đầu một cách khó nhọc, cử động gượng gạo như bị gỉ sét - ``chỉ có đầu chị là còn cử động được thôi.''

Nàng Y tá giật mình với chiếc đầu quay đi quay lại giật giật, lùi lại một chút: ``Sợ thế!''

Roboco lại quay đầu trở về vị trí cũ, gần như không thể kiểm soát được tốc độ chuyển động. ``Giờ chị phải tìm cách sửa thôi. Chị không thể live-stream như thế này được\dots''

Trên đầu nàng tóc vàng lơ lửng những dấu hỏi. ``Em thì không biết gì về máy móc cả\dots\ Giá mà Kuon-sama ở đây thì tốt quá\dots''

``Ừ\dots\ hôm nay anh ấy lại có việc ở nhà, thế có chán không chứ,'' Roboco tỏ vẻ tiếc nuối. ``Cứ thế này chắc chị cũng chẳng thể cử động được đấy chứ.''

Trong lúc Choco vẫn đang rối bời không biết phải làm gì với tình trạng của tiền bối cô, một giọng nói quen thuộc vang lên từ ngoài cửa.

``Chào mọi người-nanodesu!'' Nàng Cô đồng lấp ló ngó vào, đôi mắt đỏ đảo quanh căn phòng một cách tò mò.

Nhưng ngay khi ánh nhìn của cô chạm đến Roboco, nàng Người máy lập tức rên lên như một con thú bị dồn vào đường cùng. Cô quay ngoắt đầu về phía Rushia, phát ra một câu tiếng Anh bập bẹ với giọng méo mó: ``He-rư-pư mi\dots! (\eng{Cứu chị\dots!})''

\img{7-2b}

Choco chỉ biết đứng đó, nhìn nàng Người máy với vẻ ái ngại, trong khi Rushia thì cứng đờ, đầu óc chưa kịp xử lý được cảnh tượng kỳ lạ trước mắt.

``\textit{Có kẻ kỳ quặc trong phòng kìa!}'' Cô đồng hét lên trong lòng, không dám thốt ra thành tiếng.

Nhận thấy hậu bối của mình có vẻ đang sợ hãi, Roboco vội vàng lên tiếng giải thích, vẫn giữ nguyên giọng điệu máy móc khi nãy: ``Chị lười đi bảo trì quá nên giờ chị không cử động được.''

Rushia lập tức rùng mình. ``Tệ thật đó-nanodesu! Vậy hai chị đã thử gọi cho ai ở bên bảo trì chưa?''

``Chị không thể gọi điện thoại được đâu,'' Roboco giải thích. ``Chị chỉ có thể nói được thôi.''

Choco bỗng nảy ra một ý tưởng: ``Em có kinh nghiệm gì về máy móc không?''

``Không hẳn ạ,'' Cô đồng nhún vai, cười trừ một cách vô tư. ``Nhưng em biết rõ xác chết trong lòng bàn tay á!''

Vừa nói xong, cô run rẩy cái đầu như một con thây ma đang bò dậy từ lòng đất. Một tiếng *RẮC!* vang lên như tiếng xương bị gãy. Cảnh tượng đó làm nàng Rô-bốt nổi cả da gà (mà thật ra là do lỗi hệ thống): ``Eo ơi.''

\gif{7-2c}{1}{77}

Choco thấy đàn chị của mình không trông mong gì từ nàng Pháp sư, cô nổi giọng đầy tự hào: ``Máy móc có thể không phải là lĩnh vực của em, nhưng em biết rất nhiều về giải phẫu đấy!''

Nghe đến đó, Roboco bắt đầu đổ mồ hôi hột. ``\textit{Chẳng lẽ cô ấy định giải phẫu mình luôn\dots!?}''

``Em sẽ tiến hành bảo trì cho chị, Roboco-sama!'' nàng Y tá nở một nụ cười rạng rỡ, móc ra từ trong túi một cái cờ lê sáng bóng, đôi mắt sáng lóa lên.

\img{7-2d}

Roboco lập tức rùng mình lùi lại theo phản xạ, hay chính xác hơn, cô muốn lùi nhưng cơ thể không cho phép: ``Nhưng mà em nói là em không biết gì về máy móc mà!''

Choco nghiêng đầu, tay tung hứng cái cờ lê, giọng toát lên một sự tự tin cao đô: ``Chị có hình dáng giống con người, nên chắc là sẽ ổn thôi.''

``\textit{Không phải vậy đâu!}'' nàng Người máy hét lên trong lòng, nhưng không thể phản kháng được với tình trạng hiện tại.

``Đừng lo, Roboco-sama!'' Choco cười tươi, bắt đầu tiến lại gần. ``Em sẽ giúp chị!''

``Dừng lại đi, Y tá Học viện ơi!'' Roboco lắc đầu nguầy nguậy, trông không khác gì một con búp bê hỏng bị ai đó vặn vẹo quá đà.

Rushia chớp mắt, rồi nhìn thẳng vào nàng Rô-bốt với vẻ thắc mắc: ``Nhưng mà đầu chị có thể cử động bình thường nhỉ?''

Vốn đã quá quen với sự hỗn loạn trong công ty, Roboco chỉ biết thở dài. Nhưng vừa lúc đó, nàng Y tá bỗng nhiên tỏ ra tủi thân, mắt long lanh như sắp khóc tới nơi.: ``Ít nhất thì chị cũng cân nhắc nó đi chứ\dots''

``Xin lỗi em nhé,'' Roboco đành xin lỗi, ánh mắt đầy áy náy khi nhìn nàng Y tá.

Rushia ở bên cạnh cô bỗng vỗ tay một cái. ``Roboco-senpai, em đã hỏi nhiều người, và may mắn là có một người cũng đã từng phá dỡ, nên em đã gọi thử người đó rồi.''

Nghe đến đây, Roboco lập tức mừng rỡ ra mặt.``Thật sao? Cảm ơn nha!''

Lần này, cô đã có hy vọng thoát khỏi tình trạng trớ trêu này. Nhưng liệu người mà Rushia gọi đến có thực sự giúp được cô, hay lại khiến mọi chuyện rối tung hơn nữa\dots?

\ct

Sau một hồi chờ đợi, người mà Rushia gọi đã tới. Từ cửa văn phòng, Noel chạy vào, khoe cơ bắp và sức mạnh của mình: ``Dun dun! Rushia gọi em nên em đến rồi đây!''

Không biết tại sao, cô nàng lại còn\dots\ tạo bản mặt ngốc nghếch vốn có của mình: lè lưỡi, khuôn mặt khẽ đung đưa.

Roboco ngạc nhiên trước sự hiện diện của cô, người mà cô còn không ngờ tới. ``Noel-chan?''

Vẫn với động tác khoe cơ đó, cô nàng nhắng nhít tiến tới nàng Người máy, vừa đi vừa tỏ vẻ phấn khích. Cô nói: ``Em cũng đã được cậu ấy giải thích rồi. Cứ để đó cho em!''

\begin{figure}[H]
  \centering
  \begin{subfigure}[t]{0.45\textwidth}
    \centering
    \animategraphics[autoplay,loop,width=8cm]
    {30}{../../../../Image/Book?/7-2e/7-2e.}{1}{14}
  \end{subfigure}
  \quad
  \begin{subfigure}[t]{0.45\textwidth}
    \centering
    \animategraphics[autoplay,loop,width=8cm]
    {30}{../../../../Image/Book?/7-2e/7-2e.}{15}{119}
  \end{subfigure}
\end{figure}

Nàng Người máy nở một nụ cười đầy hy vọng: ``Cảm ơn em nhé, Noel-chan! Nhờ cả vào em đấy!''

``Được chứ!'' Noel đáp lại. Cô rút ra một cái chùy từ bên hông, vui vẻ như một đứa trẻ được cho phép chơi với đồ chơi mới.

\dots Im lặng.

Một sự im lặng đáng sợ bao trùm căn phòng. Không ai nói gì. Không ai nhúc nhích. Chỉ có Noel vẫn đứng đó, nhẹ nhàng đung đưa cái hông theo nhịp, tay cầm cây chùy vững chãi như thể sẵn sàng hành động ngay lập tức. Nếu có nhạc nền của H*lltaker, chắc chắn nó sẽ vừa khớp với cảnh tượng này một cách hoàn hảo\footnote{Fun fact: thực tế là có \href{https://www.youtube.com/watch?v=BWVsx8rlsnk}{một kênh chuyên dịch các video sang tiếng Việt của \hll\ đã làm một video như vậy}, và bất ngờ thay, nó lại khớp với nhịp bài nhạc một cách khá sát.}.

\gif{7-2f}{1}{48}

Roboco nuốt nước bọt, giọng đầy lo lắng. ``Ờ, Noel\dots\ Cái chùy đó thì làm được gì?''

Nàng Hiệp sĩ nở nụ cười rạng rỡ, hồn nhiên đáp lại: ``Thì nó là cái chùy đó!''

Ba người còn lại lập tức tái mặt. Không ai bảo ai, nhưng tất cả đều hiểu ra ý của cô. Roboco cảm thấy hệ thống cảnh báo của mình đang kêu inh ỏi, cô vội vã lắc đầu nguầy nguậy, lần này không phải vì bị hỏng mà là vì hoảng sợ thật sự.

\img{7-2g}

Nàng Hiệp sĩ vẫn tràn đầy tự tin, giơ cao cây chùy như thể đang khoe một món vũ khí tối tân ``Nó có thể phá dỡ máy móc một cách dễ dàng đấy!''

\gif{7-2h}{1}{100}

Roboco thét lên trong hoảng loạn: ``\textbf{Không phải phá dỡ kiểu đó!}''

Nhưng nàng Hiệp sĩ dường như không nghe thấy, vẫn vô tư vung vẩy cây chùy lên không khí, tập luyện một cách nhịp nhàng. Mỗi lần cô vung xuống, không khí xung quanh như rung lên bởi lực tác động.

Trong lúc đó, Choco ghé sát vào tai Rushia, thì thầm đầy nghi hoặc: ``Noel-sama đang hiểu nhầm gì đó sao?''

Rushia gãi đầu. ``Có lẽ `phá dỡ' không phải là từ chuẩn nhất mà em muốn nói.''

Ngay khi nghe câu đó, nàng Y tá lập tức xanh mặt. Cô liếc nhìn cô Hiệp sĩ đang vui vẻ hăng say vung chùy, rồi nhìn lại Rushia, giọng điệu đầy châm chọc nhưng không giấu nổi sự lo lắng: ``Đúng là đầu óc ngu si mà tứ chi thì phát triển!''

Khi đã khởi động xong xuôi, Noel nắm chặt cây chùy trong tay, mặt đầy tự tin, từ từ tiến tới ``nạn nhân'' tội nghiệp kia.

Roboco hoảng loạn đến tột độ, lắc đầu hết mức có thể: ``\textbf{Không! Không! Dừng lại đi! Đừng làm đau chị mà!!}''

Nàng Hiệp sĩ vẫn giữ nụ cười tươi rói, đầy quyết tâm: ``Không sao đâu chị ơi! Một liệu pháp sốc sẽ giúp chị đấy! Cứ tin vào em đi!''

\gif{7-2i}{1}{111}

Roboco lúc này như muốn khóc. Cô nàng réo lên tiếng cảnh báo vang khắp cả văn phòng: ``Cứu với! Ai đó làm ơn cứu tôi đi!!''

Trong khi Noel hớn hở tiến tới, Roboco hoảng loạn đến mức suýt khóc, còn Rushia đứng đơ tại chỗ, Choco nhận ra tình hình không thể để tiếp diễn như vậy được nữa.

``Được rồi, Noel-sama, dừng lại ngay!'' nàng Y tá nhanh tay ôm chặt lấy cô Hiệp sĩ từ phía sau, cố gắng giữ cô lại. Nhưng vấn đề là Noel không phải người có thể dễ dàng bị cản trở như vậy.

``Ế? Cô làm gì thế, Choco-sensei?'' Noel vẫn tiếp tục tiến tới và cầm cái chùy, mặc kệ việc có một người đang cố hết sức giữ mình lại. Choco nghiến răng, cố hết sức kéo lùi cô ấy lại, nhưng lực của một Hiệp sĩ vạm vỡ cũng không phải dạng vừa.

``Trời ơi\dots\ Noel-sama này khỏe quá\dots! Cô không thể cản em ấy được!'' Nhưng càng giữ, Noel càng tiến tới, như một chiếc xe tăng từ từ lăn bánh. Roboco nhìn cảnh tượng trước mắt mà phát hoảng, cơ thể bỗng run rẩy dữ dội.

``\textbf{Rushia-sama! Gọi Kuon-sama ngay!!}'' Choco hét lên.

Rushia như bừng tỉnh khỏi cơn sốc, vội vàng rút điện thoại ra, run rẩy bấm số. Trong lòng cô biết rất rõ, Kuon đang có việc riêng ở nhà, nhưng tình hình bây giờ không còn lựa chọn nào khác.

``Làm ơn, bắt máy đi, bắt máy đi\dots''

Cuối cùng, đầu dây bên kia cũng có tiếng trả lời. ``A-lô?''

Ngay khi cậu Tổng bắt máy, nàng Pháp sư hối hả nói: ``Kuon-senpai! Giúp bọn em với!''

``Lại cái chuyện gì ngớ ngẩn nữa hả?'' cậu đáp lại. ``Mấy đứa đúng là không thể yên ổn được mà. Sao? Nói đi. Nhà anh đang có việc.''

``Noel đang định phá Roboco-senpai kìa! Giúp với!''

``Em bị làm sao đấy? Anh tưởng là PON-chan khỏe lắm mà?'' cô cảm thấy cậu Tổng đang nhướng mày, mặc dù nàng Pháp sư không nhìn thấy được.

``Nhưng mà chị ấy đang bị đơ với hỏng rồi kìa!'' Nói đoạn, cô hướng chiếc điện thoại về phía ba cô nàng: một người đang giữ một người siêu khỏe, một người siêu khỏe thì vẫn đang tiến tới người đang hỏng hóc, còn người hỏng hóc thì lẩy bẩy cả người mà không thể làm gì được.

Bên kia, cậu con trai thở dài, giọng đầy mệt mỏi. Cậu chẹp miệng: ``Giời ạ\dots\ Lại cái tội không bảo trì đây mà\dots\ Giữ em ấy lại đi, để anh chạy tới giúp.''

Sau đó, Rushia nghe thấy tiếng vọng lại từ đầu dây bên kia, cũng từ chính Kuon, nhưng lần này là một tiếng chửi thể trong sự bực tức mà tiếng mẹ đẻ của cậu: ``\vie{Đệch mợ khổ lắm cơ, đến cả dự đám ma nhà cũng đéo yên nữa\dots}''

Cô nàng không hiểu cậu vừa chửi cái gì, nhưng cái giọng điệu kia\dots\ nghe thôi cũng đủ khiến cô toát mồ hôi hột. ``Chắc anh ấy bực lắm đây\dots''

Bây giờ chỉ còn hy vọng Kuon có thể tới kịp trước khi Roboco bị ``phá dỡ'' bởi sức mạnh hủy diệt của Noel\dots

\ct

May mắn thay, hỉ sau một hồi, Kuon đã kịp lao tới văn phòng, trên đầu vẫn còn đeo chiếc khăn tang trắng. Dù còn đang bận chuyện nhà, cậu vẫn phải tạm gác lại để xử lý đống hỗn độn này.

Với những thao tác thuần thục, Kuon tháo vát sửa lại cho Roboco, vừa làm vừa lẩm bẩm đầy khó chịu. Trong lúc cô nàng đang được kiểm tra lại lần cuối, Roboco vừa gãi đầu vừa cười trừ: ``Hầy, nãy giờ sợ thật\dots\ Em không biết phải làm gì luôn\dots''

``Đấy, cái tội lười bảo trì đấy, cái đại PON này. Ai bảo mải stream chơi game cho lắm vào,'' cậu Tổng nhăn mặt, tỏ vẻ rất khó chịu.

``Em không có PON mà\~{}!! Em chỉ là quên bảo trì thôi!'' nàng Người máy phản pháo, cảm thấy mệt mỏi vì bị cậu trêu quá nhiều. Mặc dù lần này cậu không có hàm ý chê một chút nào.

Cậu con trai nhíu mày, nhìn cô nàng Rô-bốt với ánh mắt đầy vẻ ``cạn lời'': ``Thế sao em không quên mịa luôn cả cái thân ở nhà luôn đi!?''

Thấy cậu nổi nóng một cách bất thường, Choco vội vàng tìm cách để hòa hoãn: ``Thôi nào, Kuon-sama. Đừng có mà ác ý với chị ấy chứ. Bọn em cũng đâu có ngờ tới\dots''

Cậu quay sang Rushia, gương mặt bực bội hơn hẳn: ``Mà cả Rushia-chan nữa, nói thế quái nào mà Noel-chan lại vác cái chùy đến thế kia?''

``Thì em chỉ nói là `phá dỡ' chứ em đâu có ngờ là như vậy đâu!'' nàng Pháp sư đành phải tự bào chữa.

``Sao không nói mịa luôn từ `sửa chữa' ấy?'' Kuon đáp lại, cao giọng. ``Bao nhiêu người mấy đứa không gọi mà sao cứ nhắm vào anh? Cái đám nhân viên sửa chữa máy móc của công ty vứt vào sọt rác hết à!?''

``Bọn em không biết số! Với cả bọn em bí quá, chỉ có thể nghĩ tới việc gọi cho anh thôi-nanodesu!''

Kuon đưa tay ôm đầu, thở dài một hơi. ``Anh bó tay với mấy đứa luôn. Anh không đến văn phòng một ngày là lại xảy ra chuyện ngay\dots''

``Thôi, dù gì Roboco-sama cũng cử động bình thường được rồi,'' nàng Y tá nói. ``Cảm ơn anh đã giúp chị ấy, Kuon-sama.''

``Không có gì đâu,'' cậu Tổng đáp lại. ``Nhưng mà, lần sau đừng để tình hình trở nên như thế nữa. Anh không thể đến giúp mọi lúc được đâu. Nhất là mấy ngày tới anh còn nghỉ Tết nữa, mấy đứa sẽ xử lý thế nào?''

``Cái đó thì\dots\'' Rushia ngập ngừng, không tìm được câu trả lời. Cô lập tức quay ngoắt sang chủ đề khác: ``Dù sao thì, em cũng mừng cho chị đấy,'' nàng Pháp sư nói, khẽ cúi đầu trước cậu.

``Vậy là xong rồi chứ gì? Thôi, anh về nhà lo việc tang tiếp đây.'' Cậu vừa quay lưng đi vừa lầm bầm, giọng đầy bực tức: ``\vie{Mẹ, nhà bao việc, điên thế\dots}''

Không khí trong phòng bỗng nhiên lạnh đi mấy phần. Roboco, Choco, và Rushia thoáng run rẩy khi cảm nhận cơn giận bị kiềm nén của cậu quản lý. Khi đi ngang qua Roboco, cậu liếc cô nàng một cái sắc lẹm, như muốn nhắc nhở rằng chính cô là nguyên nhân khiến cậu phải rời nhà đến đây.

Ngay khi Kuon vừa bước khỏi phòng, bầu không khí căng thẳng cũng dần tan biến\dots

\dots cho đến khi Noel hớn hở quay lại, như thể chưa từng có gì xảy ra. ``Dun dun! Vậy là sửa xong rồi hả?''

Roboco giật bắn mình, mặt tái mét, suýt lại đứng đơ lần nữa.

\img{7-2j}

Nàng Y tá và Cô đồng nhìn nhau, trên trán đẫm mồ hôi lạnh. ``\textit{Roboco hình thành nỗi sợ Noel từ đó à\dots}''

Có lẽ sẽ mất khoảng thời gian để có thể giúp cô nàng Người máy hết sợ nàng Hiệp sĩ đô con nhưng ``não nho'' này.
%==========================================TRUYỆN PHỤ=========================================
\omk

Ngay khi trông thấy Noel, Roboco gần như phát hoảng và lao đi tán loạn khắp phòng như một con gà mất đầu. Chẳng ai có thể hiểu được phản xạ mạnh mẽ này của cô nàng Rô-bốt, nhưng có vẻ như nỗi sợ hãi từ lúc trước vẫn còn ám ảnh quá nặng.

``Á! Noel khỏe quá\dots''  Rushia chật vật giữ chặt lấy cô nàng Hiệp sĩ, nhưng chẳng khác nào cố ghì một con bò tót đang hăng máu.

``Giữ chặt em ấy vào! Để chị thử gọi cho anh ấy lần nữa\dots'' nàng Y tá hốt hoảng cố tìm chiếc điện thoại, nhưng bị nàng Pháp sư ngăn cản: ``Không được! Chị vừa thấy Kuon-senpai nổi giận thế nào mà!''

Vừa nhắc đến cậu con trai ấy thôi, cả hai đã bất giác rùng mình nhớ lại ánh mắt đáng sợ của cậu trước khi rời đi.

``Nhưng nếu không thì ai sẽ giúp chúng ta bây giờ?'' Choco vừa tìm cách giữ một trong hai người, vừa hỏi.

``\dots hỉ còn cách tự xử thôi- Á! Noel, bình tĩnh lại đi!'' Rushia đáp lại, cố gắng ghìm chặt nàng Hiệp sĩ trong vô vọng.

``Hai người bị sao vậy? Thả tớ ra!'' Dù không hiểu chuyện gì đang xảy ra, cô nàng Hiệp sĩ vẫn quẫy đạp quyết liệt, cố gắng vùng khỏi sự kìm kẹp của hai người kia.

Cùng lúc đó, Roboco, trong cơn hoảng loạn tột độ, vẫn tiếp tục chạy loạn khắp phòng, trông chẳng khác nào một con chuột bị dồn đến đường cùng. Nhưng rồi, khi quá túng quẫn, nàng rô-bốt bỗng cắm đầu lao thẳng về phía cửa sổ.

\img{7-2k}

``\textbf{K-Khoan đã!!}''

*RÀM!* *XOẢNG*

Chẳng ai kịp ngăn cản, nàng Người máy đã phá tung cửa sổ và nhảy xuống đường, rồi tiếp tục chạy như thể muốn bỏ trốn khỏi thế giới này.

Tình hình còn tệ hơn khi Noel, vừa thấy Roboco bỏ chạy, lập tức hứng khởi đuổi theo. ``Ơ? Chờ đã, Roboco-senpai, đợi em với!'' Cô nàng Hiệp sĩ vui vẻ lao ra ngoài, mắt sáng rực như trẻ con nhìn thấy trò chơi mới.

Hai người còn lại chỉ biết đứng tại chỗ, thở dài đầy bất lực.

\begin{dialogue}
  \speak{Rushia} \dots Thế là Noel cũng chạy theo luôn-nanodesu.
  \speak{Choco} Còn cách nào khác đâu\dots\ Đi theo bọn họ thôi.
\end{dialogue}

Và thế là, cả bốn người bắt đầu một cuộc rượt đuổi náo loạn khắp thành phố. Roboco chạy như điên, Noel vô tư đuổi theo, còn Rushia và Choco thì hì hục bám sát phía sau, chỉ mong sao có thể kết thúc chuyện này trước khi nó vượt quá tầm kiểm soát.

Cuộc rượt đuổi kéo dài đến tận xế chiều, và cuối cùng, công ty phải điều động hẳn một nhóm nhân viên để bắt hai người kia lại.

Sau khi Roboco bị giữ lại, nàng Rô-bốt tội nghiệp vẫn còn run rẩy, bám chặt lấy cánh tay của Choco như thể vừa sống sót qua một cơn ác mộng. ``C-C-Chị sợ lắm hai đứa ơi\dots\ Ai đó ngăn cản\dots\ Noel-chan\dots\ giúp chị với\dots''

Cả nàng Y tá và nàng Pháp sư nhìn nhau lần nữa, thở dài vì mệt mỏi và chán chường sau những sự kiện trong ngày.

\begin{dialogue}
  \speak{Rushia} Quả không có Kuon-senpai ở đây thì cũng khó nhỉ\dots
  \speak{Choco} Cô phải đồng ý với em\dots *thở dài*
\end{dialogue}

Dù vậy, mọi chuyện cũng đã tạm ổn.

Nhưng không may cho Roboco, nỗi sợ Noel (hay Noel-phobia) của cô không biến mất ngay lập tức. Trong suốt vài ngày sau đó, hễ thấy bóng dáng nàng Hiệp sĩ lởn vởn đâu đó là cô lại giật bắn mình, trốn biệt vào một góc.

Cuối cùng, công ty phải nhờ đến Kuon để giúp Roboco điều trị tâm lý. Dĩ nhiên, cậu Tổng cũng chẳng mấy vui vẻ khi lại phải giải quyết một vấn đề khác do sự ``PON'' của ai đó gây ra, nhưng ít ra thì cậu không còn nổi nóng như ngày hôm đó.

``Này, đến cái nỗi sợ vớ vẩn này mà em cũng cần anh giúp nữa hả?'' cậu nhướng mày, nói với tông giọng mỉa mai. ``Với cả, anh cũng hơi ngạc nhiên là robot cũng biết sợ đấy.''

``Em\dots\ em không kiểm soát được\dots\ Chỉ cần thấy Noel thôi là em lại hoảng!'' nàng Người máy run lên như cầy sấy, thi thoảng còn giật mình như thể bị mê sảng.

``Thế thì tập đối diện đi. Có phải trẻ con đâu mà cứ né tránh mãi được.''

Và thế là, dưới sự hướng dẫn của cậu Tổng, nàng Rô-bốt đành phải trải qua một loạt bài tập tiếp xúc với Noel, từ đứng gần, nói chuyện, cho đến\dots\ chịu đựng cô nàng ôm chặt với vẻ mặt vui tươi thường thấy. Tất nhiên là đã có sự giám sát của cậu và mọi người.

Phải mất mấy ngày liền, Roboco mới có thể bình tĩnh lại khi gặp Noel mà không chạy mất dạng.

Đó là cách mà Noel-phobia của cô nàng được chữa trị thành công. Dù vậy, từ đó về sau, hễ thấy Noel vung chùy tập luyện, Roboco vẫn sẽ lùi lại một chút\dots\ chỉ để đảm bảo an toàn.

\end{document}